% problem-000331 3 阿司匹林两步滴定
\section*{3. 阿司匹林两步滴定}
虽然⼄酰⽔杨酸（相对分⼦质量180.2）含有羧基可⽤NaOH溶液直接滴定，但⽚剂中往往加⼈少量酒⽯酸或柠檬酸稳定剂，制剂⼯艺过程中也可能产⽣⽔解产物（⽔杨酸、醋酸），因此宜采⽤两步滴定法。即先⽤NaOH溶液滴定样品中共存的酸（此时⼄酰⽔杨酸也⽣成钠盐），然后加⼊过量NaOH溶液使⼄酰⽔杨酸钠在碱性条件下定量⽔解，再⽤ $\mathrm { H _ { 2 } S O _ { 4 } }$ 溶液返滴定过量碱。

（图：）
⼄酰⽔杨酸

\subsection*{3-1}
写出定量⽔解和滴定过程的反应⽅程式。

\subsection*{3-2}
定量⽔解后，若⽔解产物和过量碱浓度均约为 $0 . 1 0 \mathrm { m o l } \mathrm { L } ^ { - 1 }$ $\mathrm { H _ { 2 } S O _ { 4 } }$ 溶液浓度为 $0 . 1 0 0 0 \mathrm { m o l } \mathrm { L } ^ { - 1 }$ ，则第⼆步滴定的化学计量点pH是多少？应选⽤什么指⽰剂？滴定终点颜⾊是什么？若以甲基红作指⽰剂，滴定终点pH为4.4，则⾄少会有多⼤滴定误差？（已知⽔杨酸的离解常数为 $\mathrm { p } K _ { a 1 } = 3 . 0 , \mathrm { p } K _ { a 2 } = 1 3 . 1$ ；醋酸的离解常数为 ${ \sf N } _ { a } = 1 . 8 { \sf \times } 1 0 ^ { - 5 } )$ ）。

\subsection*{3-3}
取10⽚复⽅阿司匹林⽚，质量为� (g)。研细后准确称取粉末 $\mathbf { \dot { \tau } } m _ { 1 } \mathbf { \tau } ( \mathbf { g } )$ ，加20mL⽔（加少量⼄醇助剂），振荡溶解，⽤浓度为� (mol L−1)的NaOH溶液滴定，消耗 $: V _ { 1 } ~ ( \mathrm { m L } )$ ，然后加⼈浓度为 $\lvert c _ { 1 }$ 的 NaOH 溶液 $V _ { 2 }$ $( \mathrm { m L } )$ ；再⽤浓度 $c _ { 2 } ( \mathrm { m o l } \mathrm { L } ^ { - 1 } )$ 的 $\mathrm { H _ { 2 } S O _ { 4 } }$ 溶液滴定，消耗 $V _ { 3 } ~ \mathrm { \Omega } ( \mathrm { m L } )$ ，试计算⽚剂中⼄酰⽔杨酸的含量（mg/⽚）。
